\section{Explicit Fermat-twelve isogeny and bounded Hodge-cycle
suppliers}\label{explicit-fermat-twelve-isogeny-and-bounded-hodge-cycle-suppliers}

Author/checker: GPT-6 Astra Pro; inherited effort not exposed;
HODGE\_DIVISION /root/sixfolds. Date: 2026-11-21. New quantitative
proof; independent whole-proof review PASS reported by root
on2026-11-21, with the same model and inherited-effort metadata.

Write \(C=F_{12}\), \(J=J(C)\), and let \(\Theta\) be its canonical
principal polarization. Put \(N=144\). The explicit conclusions are

\[
 \phi:E_i^{27}\times E_\omega^{28}\longrightarrow J,\qquad
 \deg\phi\le144^{55},
 \tag{F12.1}
\]

and, after choosing the elliptic coordinates described below,

\[
 \frac14\Theta_0\ \le\ \phi^*\Theta\
       \le4\cdot144^2\Theta_0
 \tag{F12.2}
\]

in the cone of nef rational divisor classes. Here \(\Theta_0\) is the
product principal polarization.

For every \(s\ge1\), put \(g_s=55s\) and let
\(\Theta_s=\sum_{a=1}^s\operatorname{pr}_a^*\Theta\) on \(J^s\). There
is an explicitly indexed finite set of effective integral cycles of
codimension \(k\), for each \(0\le k\le g_s\), whose classes span

\[
 H^{2k}(J^s,\mathbf Q)\cap H^{k,k}(J^s)
\]

and whose degrees satisfy

\[
 \deg_{\Theta_s}Z
 \le 8^k(55s)!\,144^{55s}.
 \tag{F12.3}
\]

The statement is rational spanning by actual effective integral cycles.
It does not assert that these classes form a saturated basis of the
integral Hodge lattice. Such spanning is sufficient to choose a
bounded-degree cycle pairing nontrivially with any specified nonzero
rational Hodge class of complementary degree.

\subsection{SF-F1201. Integral projectors and the exact elliptic
census}\label{SF-F1201}

Use the affine equation \(x^{12}+y^{12}=1\). The diagonal automorphisms

\[
 \sigma_1(x,y)=(\zeta_{12}x,y),\qquad
 \sigma_2(x,y)=(x,\zeta_{12}y)
\]

generate \(G\simeq(\mathbf Z/12)^2\). Adjunction for this smooth plane
curve gives the holomorphic differential basis

\[
 \eta_{a,b}=x^{a-1}y^{b-12}\,dx,\qquad
 1\le a,b,\quad a+b<12.
 \tag{F12.4}
\]

The orders at \(x=0\) and \(y=0\) are respectively \(a-1\) and \(b-1\);
for the latter use \(dx=-y^{11}x^{-11}dy\). At infinity put \(z=1/x\)
and \(v=y/x\), so \(1+v^{12}=z^{12}\); the order is \(11-a-b\). These
orders are nonnegative. The fifty-five forms are independent because
their characters are distinct, and they exhaust the genus
\((12-1)(12-2)/2=55\). Their \(G\)-characters are \((a,b)\), meaning
\(\sigma_1^*\eta=\zeta_{12}^a\eta\) and
\(\sigma_2^*\eta=\zeta_{12}^b\eta\). The entire \(H^1(J,\mathbf C)\)
therefore has the 110 characters

\[
 \mathcal X=\{(a,b):1\le a,b\le11,\ a+b\not\equiv0\pmod{12}\},
\]

each with multiplicity one.

Let \(O\) run through the orbits of \((\mathbf Z/12)^*=\{1,5,7,11\}\) on
\(\mathcal X\), acting by simultaneous multiplication. Define the
rational idempotent

\[
 e_O=\frac1{144}\sum_{x,y\bmod12}
       \left(\sum_{(a,b)\in O}\zeta_{12}^{-ax-by}\right)
       \sigma_1^x\sigma_2^y.
 \tag{F12.5}
\]

The inner coefficient is a rational algebraic integer, hence an integer.
Thus \(144e_O\in\mathbf Z[G]\). Character orthogonality shows that these
projectors are mutually orthogonal and sum to the identity on
\(H^1(J)\). The corresponding identities hold as rational endomorphisms
of \(J\), since an endomorphism acting trivially on rational first
homology is zero.

Put \(A_O=\operatorname{im}(144e_O)\), with its connected image as an
abelian subvariety of \(J\). Its dimension is \(|O|/2\), which is one or
two. Let

\[
 O^+=\{(a,b)\in O:a+b<12\}.
\]

The following table gives a complete orbit census. A representative
stands for its orbit under the four units; the rows contain disjoint
orbits.

\begin{longtable}[]{@{}
  >{\raggedright\arraybackslash}p{(\columnwidth - 6\tabcolsep) * \real{0.2143}}
  >{\raggedleft\arraybackslash}p{(\columnwidth - 6\tabcolsep) * \real{0.2857}}
  >{\raggedleft\arraybackslash}p{(\columnwidth - 6\tabcolsep) * \real{0.2857}}
  >{\raggedright\arraybackslash}p{(\columnwidth - 6\tabcolsep) * \real{0.2143}}@{}}
\toprule\noalign{}
\begin{minipage}[b]{\linewidth}\raggedright
Character conductor
\end{minipage} & \begin{minipage}[b]{\linewidth}\raggedleft
Pure CM field
\end{minipage} & \begin{minipage}[b]{\linewidth}\raggedleft
Factor dimension
\end{minipage} & \begin{minipage}[b]{\linewidth}\raggedright
Number of factors
\end{minipage} \\
\midrule\noalign{}
\endhead
\bottomrule\noalign{}
\endlastfoot
\(12\) & \(\mathbf Q(i)\) & \(2\) & \(12\) \\
\(12\) & \(\mathbf Q(\omega)\) & \(2\) & \(9\) \\
\(6\) & \(\mathbf Q(\omega)\) & \(1\) & \(9\) \\
\(4\) & \(\mathbf Q(i)\) & \(1\) & \(3\) \\
\(3\) & \(\mathbf Q(\omega)\) & \(1\) & \(1\) \\
\end{longtable}

Complete representative lists, in the same row order, are:

\begin{itemize}
\item
  Conductor12, Gaussian: \((1,1)\), \((1,3)\), \((1,5)\), \((1,6)\),
  \((1,8)\), \((1,10)\), \((2,5)\), \((3,1)\), \((3,4)\), \((4,3)\),
  \((4,5)\), \((6,1)\).
\item
  Conductor12, Eisenstein: \((1,2)\), \((1,4)\), \((1,7)\), \((1,9)\),
  \((2,1)\), \((2,3)\), \((3,2)\), \((3,7)\), \((4,1)\).
\item
  Conductor6: \((2,2)\), \((2,4)\), \((2,6)\), \((2,8)\), \((4,2)\),
  \((4,6)\), \((4,10)\), \((6,2)\), \((6,4)\).
\item
  Conductor4: \((3,3)\), \((3,6)\), \((6,3)\).
\item
  Conductor3: \((4,4)\).
\end{itemize}

Here the conductor is \(12/\gcd(12,a,b)\). The field label has the
following explicit integral meaning. For each Gaussian row, one of
\(\sigma_1^3,\sigma_2^3\) acts by the same primitive fourth root on all
the characters of \(O^+\). For each Eisenstein row, one of
\(\sigma_1^4,\sigma_2^4\) acts by the same primitive third root on all
the characters of \(O^+\). This is verified by reducing the displayed
orbit coordinates modulo four or three, respectively. The certificate
records the chosen automorphism and scalar for every orbit.

Consequently \(A_O\) admits the full maximal order \(\mathbf Z[i]\) or
\(\mathbf Z[\omega]\) as integral endomorphisms, with a single complex
embedding on its entire tangent space. This is stronger than merely
identifying a rational CM field.

Both orders are Euclidean for the usual norm: nearest Gaussian and
Eisenstein lattice points have squared distances at most \(1/2\) and
\(1/3\), respectively. In particular both are principal ideal domains.
If \(A_O=V/\Lambda\) has dimension \(d\), its lattice \(\Lambda\) is a
torsion-free module of rank \(d\) over the indicated order, hence free.
An order basis of \(\Lambda\) is a complex basis of \(V\), because its
complex span contains the full lattice and therefore all of \(V\). It
follows that

\[
 A_O\simeq E_i^d\quad\hbox{or}\quad A_O\simeq E_\omega^d
 \tag{F12.6}
\]

as unpolarized abelian varieties. The maps are analytic homomorphisms
between projective complex tori and therefore algebraic. This argument
introduces no further isogeny degree.

The table has thirty-four factors. Summing their dimensions gives

\[
 12\cdot2+3=27,\qquad
 9\cdot2+9+1=28.
\]

This proves the asserted elliptic type by a complete finite computation
and the integral lattice argument, without a height estimate.

\subsection{SF-F1202. Polarization exponents and the isogeny
degree}\label{SF-F1202}

Every \(g\in G\) preserves the Jacobian polarization, so its Rosati
adjoint is \(g^{-1}\). Each orbit \(O\) is stable under negation. Thus
(F12.5) gives \(e_O^\dagger=e_O\). Distinct factor spaces are orthogonal
for the polarization, and the addition map

\[
 \phi:\prod_O A_O\longrightarrow J
 \tag{F12.7}
\]

is an isogeny with

\[
 \phi^*\Theta=\boxplus_O\Theta_O,\qquad
 \Theta_O=\Theta|_{A_O}.
 \tag{F12.8}
\]

There is also an integral reverse map \(J\to\prod_O A_O\), with
coordinates \(144e_O\); both composites are multiplication by 144. This
already gives a degree bound, but the polarizations give the sharper
bound used here.

For a fixed factor \(A=A_O\), its polarization complement \(A^\perp\) is
the connected kernel of the surjection \(J\to\widehat A\) obtained from
the principal polarization and dual inclusion. Hence

\[
 \ker\lambda_{\Theta_O}=A\cap A^\perp.
\]

The endomorphism \(144e_O\) is multiplication by 144 on \(A\) and zero
on \(A^\perp\). It follows that this finite polarization kernel is
killed by 144. Every elementary divisor in the polarization type
therefore divides 144. If \(d=\dim A_O\),

\[
 \chi(\Theta_O)\le144^d.
 \tag{F12.9}
\]

Since \(\Theta\) is principal, the top-intersection formula for a finite
pullback and the product formula in (F12.8) give

\[
 \deg\phi=\chi(\phi^*\Theta)
       =\prod_O\chi(\Theta_O)\le144^{55}.
 \tag{F12.10}
\]

Combining (F12.6) and (F12.7) proves (F12.1).

\subsection{SF-F1203. A uniform comparison with the product
polarization}\label{SF-F1203}

The one-dimensional factors have degree between one and 144, so their
polarizations already lie between the product principal polarization and
144 times that polarization.

Consider a two-dimensional factor \(E_K^2\), with \(K=\mathbf Q(i)\) or
\(\mathbf Q(\omega)\). In an order basis, its polarization is
represented by a positive definite integral Hermitian matrix

\[
 H=\begin{pmatrix}a&b\\\bar b&c\end{pmatrix},
 \qquad a,c\in\mathbf Z,\quad b\in\mathcal O_K,\quad
 \mathcal D=ac-|b|^2=\chi(\Theta_O)\le144^2.
 \tag{F12.11}
\]

There is no extra factor in the last equality: relative to the product
principal polarization, the polarization homomorphism is the Hermitian
endomorphism \(H\). Its degree is
\(\operatorname{Nm}_{K/\mathbf Q}\det H=\mathcal D^2\), while the degree
of a polarization homomorphism is its Euler characteristic squared.

Choose a nonzero vector of smallest \(H\)-norm, denoting that norm by
\(a\). It is primitive over \(\mathcal O_K\), since division by a
nonunit common factor would strictly decrease its positive norm. Extend
it to an order basis. Replace the second vector by an order multiple of
the first so that

\[
 |b|^2\le\rho^2a^2,\qquad
 \rho^2=\frac12\ \text{or}\ \frac13.
\]

Minimality gives \(c\ge a\), and consequently

\[
 \mathcal D\ge(1-\rho^2)a^2\ge\frac{a^2}{2}.
\]

Since \(a\) is a positive integer,

\[
 a\le\sqrt{2\mathcal D}\le2\mathcal D,\qquad
 c=\frac{\mathcal D+|b|^2}{a}
    \le\mathcal D+\frac a2\le2\mathcal D.
 \tag{F12.12}
\]

Thus \(\operatorname{tr}H\le4\mathcal D\le4\cdot144^2\). Also

\[
 \operatorname{tr}H^{-1}
   =\frac{\operatorname{tr}H}{\mathcal D}\le4.
\]

A positive Hermitian matrix has maximum eigenvalue at most its trace. It
follows that

\[
 \frac14I\le H\le4\cdot144^2I.
 \tag{F12.13}
\]

Taking the direct sum over all thirty-four factors proves (F12.2). The
inequalities mean that the differences are represented by semipositive
translation-invariant forms; equivalently they are nef rational divisor
classes. Every coordinate elliptic curve in this choice of coordinates
has \(\Theta\)-degree at most \(2\cdot144^2\).

\subsection{SF-F1204. Finite effective generators on every
power}\label{SF-F1204}

Set \(B_s=E_i^{27s}\times E_\omega^{28s}\), let \(\Theta_{0,s}\) be its
product principal polarization, and let

\[
 \Phi_s=\phi^s:B_s\longrightarrow J^s,\qquad
 L_s=\Phi_s^*\Theta_s.
\]

Then

\[
 \Theta_{0,s}\le4L_s,\qquad
 \chi(L_s)=(\deg\phi)^s\le144^{55s}.
 \tag{F12.14}
\]

For each of the two CM blocks, take the following effective divisors:

\begin{enumerate}
\def\labelenumi{\arabic{enumi}.}
\item
  the coordinate kernels \(x_j=0\);
\item
  for every \(j<k\), the kernels \(x_j-x_k=0\);
\item
  for every \(j<k\), the kernels \(x_j-\tau_Kx_k=0\), with \(\tau_i=i\)
  and \(\tau_\omega=\omega\).
\end{enumerate}

There are

\[
 (27s)^2+(28s)^2=1513s^2
 \tag{F12.15}
\]

such divisors. Their Hermitian matrices are the coordinate rank-one
matrices and the matrices \(vv^\dagger\) with \(v\) supported on two
coordinates of unit absolute value. Thus every generator \(D_j\)
satisfies

\[
 0\le D_j\le2\Theta_{0,s}\le8L_s.
 \tag{F12.16}
\]

Their classes span the rational divisor space: the coordinate classes
supply the diagonal Hermitian entries, while the two indicated
difference kernels supply the two independent rational off-diagonal
entries in \(K=\mathbf Q+\mathbf Q\tau_K\).

Here is a direct proof that these divisors generate the entire rational
Hodge ring. Over \(F_0=\mathbf Q(i,\omega)\), decompose \(H^1(B_s)\)
according to the two embeddings of each imaginary quadratic field, and
take the resulting exterior multidegrees. A monomial has separate
imbalances \(d_i,d_\omega\), the numbers of holomorphic factors minus
antiholomorphic factors in the two blocks. Being of Hodge type \((k,k)\)
forces \(d_i+d_\omega=0\). For a rational Hodge class, all Galois
conjugate multidegree components are also Hodge components: their
projections are tensor and exterior constructions from the field-action
projectors on \(H^1\), so preserve Hodge types, and Galois conjugation
replaces each projector by its conjugate. The two quadratic fields are
distinct and linearly disjoint, so an automorphism conjugating just
\(\mathbf Q(i)\) gives \(-d_i+d_\omega=0\). Thus both imbalances are
zero.

Every balanced monomial is a product of degree-two pairings within the
individual CM blocks. Those degree-two pairings are spanned by the
divisor classes just listed. Descending the resulting spanning statement
from \(F_0\) to \(\mathbf Q\) proves the claim.

For each \(k\)-element subset of the list, form the simultaneous kernel
of its \(k\) elliptic homomorphisms. If those homomorphisms are linearly
independent in their respective CM blocks, their combined map to a
product of \(k\) elliptic curves is surjective. Its zero fibre is a
smooth, possibly disconnected, abelian kernel of codimension \(k\), and
is an effective integral cycle \(Y\) with class \(\prod D_j\). If the
maps are dependent, that product class is zero: their rank-one forms
have dependent exterior support. Discard that subset. Repeated
generators likewise give zero, so the remaining subsets span all
monomials of degree \(k\).

The finite set is explicit: test the corresponding matrices over
\(\mathbf Q(i)\) and \(\mathbf Q(\omega)\), and retain the independent
subsets. There are at most \(\binom{1513s^2}{k}\) cycles. By the
semipositive inequalities (F12.16),

\[
 \deg_{L_s}Y
 =\int_{B_s}D_{j_1}\cdots D_{j_k}L_s^{g_s-k}
 \le8^k\int_{B_s}L_s^{g_s}
 \le8^k(55s)!\,144^{55s}.
 \tag{F12.17}
\]

The pushforward \(Z=(\Phi_s)_*Y\) is effective with integral
multiplicities, and the projection formula identifies its
\(\Theta_s\)-degree with the left side. Pushforward by an isogeny is an
isomorphism on rational cohomology up to a nonzero scalar inverse to
pullback. It therefore transports the spanning statement to \(J^s\),
proving (F12.3).

If a single irreducible detector is wanted, a nonzero pairing with such
a \(Z\) implies a nonzero pairing with at least one of its irreducible
components. That component has degree at most the displayed bound. No
coefficient or saturation estimate is hidden in this last selection.

\subsection{Published source, new proof, and actual
verification}\label{published-source-new-proof-and-actual-verification}

The qualitative orbit and CM-type description is classical. A published
primary locator is Koblitz--Rohrlich, \emph{Simple Factors in the
Jacobian of a Fermat Curve}, Canadian Journal of
Mathematics30(6),1183--1205 (1978), introductory lattice decomposition
and CM-type stabilizer discussion on printed pp1183--1184, DOI
https://doi.org/10.4153/CJM-1978-099-6 . The official publisher PDF was
opened and those exact pages read. Its later theorems for prime-to-six
levels are not invoked at level twelve.

The integral projector, exponent144, degree144\^{}55, binary Hermitian
comparison, explicit finite cycle family and numerical bound above are
proved here. The exact script checks every orbit, its chosen integral
scalar action, all projector coefficients, their Fourier actions on
all144 characters, and the stated census. It does not claim to have
explicitly enumerated kernels for arbitrarily large powers or to
saturate the integral Hodge lattice.
